\documentclass[12pt,a4paper,oneside]{report}
\usepackage[a4paper,left=1.5in,right=1in,top=1in,bottom=1in]{geometry}
\usepackage{newtxtext,newtxmath}
\usepackage{setspace}
\usepackage{array,booktabs,longtable,tabularx,multirow}
\usepackage{enumitem}
\usepackage{hyperref}
\usepackage{caption}
\usepackage{graphicx}
\usepackage{float}
\usepackage{fancyhdr}
\usepackage{lastpage}
\usepackage{pifont}

\onehalfspacing
\setlength{\parindent}{0pt}
\setlength{\parskip}{6pt plus 1pt minus 1pt}
\renewcommand{\contentsname}{Contents}
\renewcommand{\bibname}{BIBLIOGRAPHY / REFERENCES}
\renewcommand{\thechapter}{\arabic{chapter}}
\renewcommand{\thesection}{\thechapter.\arabic{section}}
\renewcommand{\thesubsection}{\thesection.\arabic{subsection}}
\renewcommand{\theequation}{\thechapter.\arabic{equation}}
\captionsetup{justification=centering}
\setlength{\headheight}{28pt}
\pagestyle{fancy}
\fancyhf{}
\fancyhead[C]{\footnotesize Smart Academic Management System with AI-Based Attendance and Performance Analytics}
\fancyfoot[C]{[\thepage/\pageref{LastPage}]}
\renewcommand{\headrulewidth}{0.4pt}
\renewcommand{\footrulewidth}{0pt}
\fancypagestyle{plain}{%
  \fancyhf{}
  \fancyhead[C]{\footnotesize Smart Academic Management System with AI-Based Attendance and Performance Analytics}
  \fancyfoot[C]{[\thepage/\pageref{LastPage}]}
  \renewcommand{\headrulewidth}{0.4pt}
  \renewcommand{\footrulewidth}{0pt}
}
\let\oldsection\section
\renewcommand{\section}{\clearpage\oldsection}
\newcommand{\tick}{\ding{51}}
\newcommand{\placeholderfig}[2]{%
\begin{figure}[H]
\centering
\fbox{\begin{minipage}{0.82\textwidth}\centering\vspace{1.1cm}\textbf{Manual Image Placeholder}\par#1\par\vspace{0.35cm}\textit{Insert the image here manually.}\vspace{1.1cm}\end{minipage}}
\caption{#2}
\end{figure}}

\begin{document}

\pagenumbering{roman}
\chapter*{ABSTRACT}
\addcontentsline{toc}{chapter}{ABSTRACT}
Educational institutions increasingly require digital solutions to efficiently manage academic activities, monitor student performance, and maintain accurate records. Traditional management approaches often involve manual attendance recording, fragmented communication channels, and limited access to academic analytics, leading to inefficiencies in administrative and academic processes.

This project proposes the development of a \textbf{Smart Academic Management System with AI-Based Attendance and Performance Analytics}. The system will provide a centralized platform for managing academic activities through role-based access control for Super Admin, Department Admin, Teachers, and Students. Administrative functionalities include department, semester, course, teacher, and student management, along with institution-wide notice publication. Teachers will be able to manage attendance, assignments, internal assessments, notices, and performance evaluations. Students will have access to attendance records, academic progress statistics, assignment submissions, notices, and GPA-related insights.

A key feature of the system is the AI-powered attendance module, which utilizes YOLOv8-based face detection and facial recognition techniques to automate classroom attendance recording from captured images. Additionally, the system will provide performance analytics and GPA prediction features to support academic decision-making for both students and faculty members.

The proposed system will be developed using React.js with TypeScript for the frontend, FastAPI for backend services, PostgreSQL for data management, and WebSocket-based real-time communication. The expected outcome is a scalable, efficient, and intelligent academic management platform that improves administrative efficiency and enhances academic monitoring.

\textbf{Keywords:} Student Management System, Academic Management, YOLOv8, Facial Recognition, FastAPI, Performance Analytics, GPA Prediction, Role-Based Access Control, WebSocket, Docker.

\tableofcontents
\clearpage
\pagenumbering{arabic}

\chapter{INTRODUCTION}
\section{Project Overview}
Educational institutions manage a large volume of academic information, including student records, course allocations, attendance, assignments, notices, and performance evaluations. Traditional approaches often rely on multiple disconnected systems or manual processes, resulting in inefficiencies, data inconsistency, and increased administrative workload.

This project proposes the development of a \textbf{Smart Academic Management System with AI-Based Attendance and Performance Analytics}, a centralized web-based platform designed to streamline academic and administrative activities within an educational institution. The system will implement Role-Based Access Control (RBAC) to provide secure access to four user roles: Super Admin, Department Admin (HOD), Teacher, and Student.

The proposed system will enable academic administrators to manage departments, semesters, courses, teachers, and students efficiently. Teachers will be able to manage attendance, assignments, notices, and internal assessments, while students can access attendance records, assignment statuses, academic statistics, and GPA prediction insights.

A major feature of the project is the integration of artificial intelligence for attendance management. By using YOLOv8-based face detection and facial recognition techniques, the system will automate attendance recording from classroom images, reducing manual effort and improving accuracy. The system will also provide performance analytics and GPA prediction to help students and faculty make informed academic decisions.

\section{Problem Statement}
Many educational institutions still depend on manual or semi-digital processes for academic management. Attendance is often recorded manually, notices are communicated through scattered channels, and student performance records are maintained in separate files or spreadsheets. These practices create several problems, including data duplication, inaccurate attendance records, delayed communication, and limited visibility into student academic progress.

Manual attendance systems are time-consuming and vulnerable to proxy attendance. Teachers may spend valuable class time calling student names or checking attendance sheets. Similarly, administrators face difficulty maintaining updated records for departments, semesters, courses, teachers, and students. Students often lack a single platform where they can check attendance, assignments, notices, internal marks, and academic progress.

There is therefore a need for an integrated academic management system that combines administration, attendance automation, assignment management, real-time communication, performance analytics, and GPA prediction in a single secure platform.

\section{Project Objectives}
\textbf{General Objective:}

To develop a Smart Academic Management System that integrates academic administration, AI-based attendance tracking, and student performance analytics into a centralized platform.

\textbf{Specific Objectives:}
\begin{enumerate}
\item To implement a secure role-based access control system for Super Admin, Department Admin, Teachers, and Students.
\item To manage departments, semesters, courses, teachers, and student enrollments through a centralized administration portal.
\item To develop an AI-powered attendance system using YOLOv8-based face detection and facial recognition techniques.
\item To provide assignment management and submission tracking functionalities.
\item To automate internal marks management and academic record maintenance.
\item To generate student performance analytics and academic progress reports.
\item To provide GPA prediction and academic recommendation features for students.
\item To implement real-time notifications and notice publishing using WebSocket communication.
\item To provide an interactive dashboard for academic monitoring and decision-making.
\end{enumerate}

\section{Significance of the Study}
The proposed system aims to improve the efficiency and effectiveness of academic management processes within educational institutions. By centralizing academic information and automating routine administrative tasks, the system reduces manual workload and minimizes the possibility of errors.

The AI-based attendance module enhances attendance accuracy while reducing the time required for attendance recording. Performance analytics and GPA prediction features provide valuable insights into student progress, enabling teachers and administrators to identify academic trends and support students more effectively.

Students will benefit from transparent access to attendance records, assignment status, academic statistics, and performance predictions, allowing them to monitor their progress and make informed academic decisions.

\section{Scope and Limitations}
\textbf{Scope:}
\begin{itemize}
\item Role-based access control for Super Admin, Department Admin, Teachers, and Students.
\item Department, semester, course, teacher, and student management.
\item Course enrollment and assignment management.
\item Notice publishing and real-time notifications.
\item AI-based attendance tracking using classroom images.
\item Internal marks management and performance evaluation.
\item Student performance analytics and GPA prediction.
\item Docker-based deployment architecture.
\end{itemize}

\textbf{Limitations:}
\begin{itemize}
\item The facial recognition module may be affected by lighting conditions, image quality, and partial facial occlusions.
\item Attendance recognition accuracy depends on the quality and quantity of training data available.
\item The system will initially support local file storage rather than cloud-based storage solutions.
\item GPA prediction results are estimations and should not be considered official academic outcomes.
\end{itemize}

\chapter{LITERATURE REVIEW}
\section{Overview of Academic Management Systems}
Academic management systems have evolved significantly over the past decade, transitioning from simple record-keeping tools to comprehensive educational platforms. Subbiah et al. (2021) developed a web-based Student Database Management System that centralized attendance, academic records, library information, and administrative functions. The system successfully replaced paper-based processes and improved data accessibility.

Yu and Wang (2022) introduced a student management system incorporating an optimized K-means clustering algorithm for intelligent student classification and performance analysis. Their research demonstrated the value of data-driven academic insights for identifying student groups with similar academic characteristics.

\section{Traditional Attendance Management Approaches}
Traditional attendance management methods, including manual roll calls and sign-in sheets, have been widely criticized for inefficiency and inaccuracy. Roy et al. (2012) studied the implementation of automated attendance systems and identified key challenges including time consumption, susceptibility to proxy attendance, and difficulties in maintaining historical records.

\section{Computer Vision-Based Attendance Systems}
\subsection{Traditional Computer Vision Techniques}
Earlier automated attendance systems primarily relied on traditional computer vision techniques. Gomes et al. (2019) developed a facial recognition attendance system using HAAR Cascade classifiers for face detection and Local Binary Pattern Histogram (LBPH) algorithms for recognition. The system successfully automated attendance recording but demonstrated recognition accuracy heavily dependent on image quality and lighting conditions.

\subsection{Deep Learning Approaches}
Recent studies have adopted deep learning approaches to overcome the limitations of traditional computer vision methods. Pratama et al. (2024) developed a real-time attendance system combining YOLOv8 for face detection and FaceNet for facial recognition, achieving an accuracy of 98.7\%.

KLE Technological University (2024) introduced a YOLOv8-enhanced one-shot learning attendance system that utilizes Siamese Networks, allowing student registration with only a single reference image while achieving approximately 97\% recognition accuracy.

\subsection{YOLOv8 for Classroom Face Detection}
Ananda et al. (2024) conducted a comprehensive evaluation of YOLOv8 for detecting low-resolution student faces in real classroom settings. Their findings demonstrated that larger YOLOv8 variants significantly improve detection accuracy while maintaining real-time detection speeds suitable for classroom applications.

\placeholderfig{FIGURE 2: Comparative analysis diagram showing existing systems vs. the proposed system}{Comparative Analysis: Existing Systems vs. Proposed System}

\section{Research Gap and Proposed Contribution}
The reviewed literature indicates that existing solutions generally focus on either student management functionalities or attendance automation independently. Attendance systems primarily emphasize facial recognition accuracy, while student management systems concentrate on administrative data handling. Very few systems provide a unified platform that combines role-based academic management, AI-powered attendance tracking, assignment management, real-time notifications, performance analytics, and GPA prediction within a single architecture.

\chapter{PROPOSED METHODOLOGY}
\section{System Architecture}
The proposed system will follow a \textbf{Modular Monolith Architecture} with an API-first design. The system consists of five major layers:

\begin{table}[h!]
\centering
\caption{System Architecture Layers}
\begin{tabularx}{\textwidth}{|>{\raggedright\arraybackslash}p{0.22\textwidth}|>{\raggedright\arraybackslash}p{0.24\textwidth}|X|}
\hline
\textbf{Layer} & \textbf{Technology} & \textbf{Purpose} \\ \hline
Presentation Layer & React.js with TypeScript & User interface with role-specific dashboards \\ \hline
Application Layer & FastAPI & Business logic, authentication, request processing \\ \hline
Database Layer & PostgreSQL & Data storage and management \\ \hline
AI Attendance Layer & YOLOv8 + Face Recognition & Face detection and student identification \\ \hline
Real-Time Layer & WebSockets & Notifications and real-time updates \\ \hline
\end{tabularx}
\end{table}

\placeholderfig{System Architecture Diagram}{System Architecture Diagram}

\section{Authentication and Authorization}
The system will implement JWT-based authentication and role-based authorization. Each user will be assigned a role that controls access to dashboards, APIs, and system features. Super Admin users will manage the overall institution, Department Admin users will manage department-level resources, Teachers will manage academic activities for assigned courses, and Students will access academic records and submissions.

\section{Functional Modules}
\begin{longtable}{|p{0.08\textwidth}|p{0.28\textwidth}|p{0.52\textwidth}|}
\caption{Functional Modules}\\
\hline
\textbf{S.N.} & \textbf{Module} & \textbf{Description} \\ \hline
1 & Super Admin Module & Institution management, departments, semesters, users, notices \\ \hline
2 & Department Admin Module & Department-level course, teacher, and student management \\ \hline
3 & Teacher Module & Attendance, assignments, notices, internal marks, evaluations \\ \hline
4 & Student Module & Attendance viewing, assignment submission, notices, marks, GPA insights \\ \hline
5 & AI Attendance Module & Face detection, recognition, attendance confirmation, report generation \\ \hline
6 & Performance Analytics Module & Attendance statistics, marks analysis, progress visualization \\ \hline
7 & GPA Prediction Module & GPA estimation and academic recommendation support \\ \hline
\end{longtable}

\section{AI-Based Attendance Workflow}
The AI-based attendance module will automate attendance recording using image capture and facial recognition. The proposed workflow is as follows:
\begin{enumerate}
\item The teacher captures or uploads a classroom image for a specific course session.
\item The backend sends the image to the AI attendance service.
\item YOLOv8 detects student faces from the classroom image.
\item Detected faces are cropped and passed to the facial recognition model.
\item Recognized student identities are matched with enrolled students for the course.
\item Attendance records are generated automatically and displayed to the teacher for verification.
\item The teacher confirms or edits the generated attendance record.
\item Final attendance data is stored in the PostgreSQL database.
\end{enumerate}

\placeholderfig{AI-Based Attendance Workflow Diagram}{AI-Based Attendance Workflow}

\section{Database Design}
The system database will be designed using PostgreSQL. Major entities include users, roles, departments, semesters, courses, enrollments, attendance, assignments, submissions, internal marks, notices, notifications, and GPA prediction records.

\begin{table}[h!]
\centering
\caption{Core Database Entities}
\begin{tabularx}{\textwidth}{|p{0.26\textwidth}|X|}
\hline
\textbf{Entity} & \textbf{Purpose} \\ \hline
Users & Stores login and profile information for all system users \\ \hline
Roles & Defines Super Admin, Department Admin, Teacher, and Student access levels \\ \hline
Departments & Stores academic department information \\ \hline
Courses & Stores course details and semester associations \\ \hline
Enrollments & Maps students to courses and semesters \\ \hline
Attendance & Stores attendance records for each course session \\ \hline
Assignments & Stores assignment details created by teachers \\ \hline
Submissions & Stores student assignment submissions and grading status \\ \hline
Internal Marks & Maintains assessment marks and evaluation records \\ \hline
Notices & Stores notices published by admins and teachers \\ \hline
\end{tabularx}
\end{table}

\section{Real-Time Communication}
WebSocket communication will be used to deliver real-time notices and notifications. When an admin or teacher publishes a notice, the backend will push the update to connected users according to their role, department, course, or enrollment group.

\section{Deployment Strategy}
The project will use Docker-based deployment. The frontend, backend, database, and supporting services will be configured through Dockerfiles and a docker-compose configuration. Nginx will be used as a reverse proxy for routing requests to frontend and backend services.

\chapter{PROPOSED PERFORMANCE ANALYSIS METHODOLOGY AND VALIDATION SCHEME}
\section{Performance Analysis Methodology}
\subsection{System Performance Metrics}
System performance will be evaluated using response time, request throughput, database query efficiency, dashboard loading time, API error rate, and concurrent user handling capacity.

\subsection{AI Attendance Performance Metrics}
The AI attendance module will be evaluated using face detection accuracy, recognition accuracy, false acceptance rate, false rejection rate, processing time per image, and attendance generation time.

\section{Validation Scheme}
\subsection{Functional Validation}
Functional validation will verify that each module performs according to the specified requirements. Test cases will cover login, RBAC access, department management, course management, attendance recording, assignment submission, notice publishing, and report generation.

\subsection{Integration Validation}
Integration validation will ensure that frontend, backend, database, AI attendance service, and WebSocket notification features work together correctly. API responses, database updates, and real-time events will be tested across user roles.

\subsection{User Acceptance Testing}
User acceptance testing will involve administrators, teachers, and students. Feedback will be collected regarding usability, correctness, response time, and overall usefulness of the system.

\chapter{PROPOSED DELIVERABLES/OUTPUT}
\section{Software Deliverables}
\begin{longtable}{|p{0.08\textwidth}|p{0.32\textwidth}|p{0.46\textwidth}|}
\caption{Software Deliverables}\\
\hline
\textbf{S.N.} & \textbf{Deliverable} & \textbf{Description} \\ \hline
1 & Admin Management Module & Department, semester, course, teacher, and student management \\ \hline
2 & Teacher Module & Attendance, assignments, marks, notices, and evaluations \\ \hline
3 & Student Module & Attendance records, assignment submissions, notices, marks, and analytics \\ \hline
4 & AI Attendance Module & YOLOv8-based face detection and recognition attendance automation \\ \hline
5 & Performance Analytics Module & Statistics, reports, and visualizations \\ \hline
6 & GPA Prediction Module & GPA estimation and academic guidance \\ \hline
7 & Notice and Notification Module & Real-time WebSocket communication \\ \hline
\end{longtable}

\section{Technical Deliverables}
\begin{itemize}
\item React.js Frontend Application (TypeScript-based)
\item FastAPI Backend Services (Python)
\item PostgreSQL Database Schema (SQL scripts)
\item AI Attendance Model Integration (YOLOv8 weights and inference code)
\item RESTful API Documentation (Swagger/OpenAPI)
\item Docker Configuration Files (docker-compose.yml, Dockerfiles)
\item Nginx Deployment Configuration
\item Source Code Repository (Git)
\end{itemize}


\chapter{PROJECT TASK AND TIME SCHEDULE}
\begin{table}[H]
\centering
\caption{Project Gantt Chart (10 Weeks)}
\label{tab:schedule}
\small
\resizebox{\textwidth}{!}{%
\begin{tabular}{|l|c|c|c|c|c|}
\hline
\textbf{Task} & \textbf{W1--2} & \textbf{W3--4} & \textbf{W5--6} & \textbf{W7--8} & \textbf{W9--10} \\
\hline
Requirement Analysis \& Lit Review       & $\checkmark$ & $\checkmark$ & & & \\
\hline
System Design \& Database Design         & $\checkmark$ & $\checkmark$ & & & \\
\hline
UI/UX Design                             & $\checkmark$ & $\checkmark$ & & & \\
\hline
Authentication \& RBAC Module            & & $\checkmark$ & $\checkmark$ & & \\
\hline
Admin Management Module                  & & $\checkmark$ & $\checkmark$ & & \\
\hline
Teacher \& Student Modules               & & & $\checkmark$ & $\checkmark$ & \\
\hline
Assignment \& Notice Management          & & & $\checkmark$ & $\checkmark$ & \\
\hline
AI Attendance Module                     & & & & $\checkmark$ & $\checkmark$ \\
\hline
Performance Analytics \& GPA             & & & & $\checkmark$ & $\checkmark$ \\
\hline
Real-Time Notification Integration       & & & & & $\checkmark$ \\
\hline
System Integration                       & & & & & $\checkmark$ \\
\hline
Testing and Validation                   & & & & & $\checkmark$ \\
\hline
Deployment (Docker \& Nginx)             & & & & & $\checkmark$ \\
\hline
Documentation                            & $\checkmark$ & $\checkmark$ & $\checkmark$ & $\checkmark$ & $\checkmark$ \\
\hline
Final Review \& Presentation             & & & & & $\checkmark$ \\
\hline
\end{tabular}%
}
\end{table}

\placeholderfig{FIGURE: Insert Gantt chart visualization manually here}{Project Gantt Chart}

\section{Milestones}
\begin{table}[H]
\centering
\caption{Project Milestones}
\label{tab:milestones}
\begin{tabular}{|c|c|p{8cm}|}
\hline
\textbf{Milestone} & \textbf{Week} & \textbf{Deliverable} \\
\hline
Milestone 1 & Week 2  & Requirements finalized, literature review completed \\
\hline
Milestone 2 & Week 4  & System architecture, database schema, UI prototypes \\
\hline
Milestone 3 & Week 6  & Authentication and core management modules \\
\hline
Milestone 4 & Week 8  & Teacher/Student modules and AI attendance module begun \\
\hline
Milestone 5 & Week 9  & AI attendance, analytics, integration, and notifications complete \\
\hline
Milestone 6 & Week 10 & Testing, deployment, documentation, final presentation \\
\hline
\end{tabular}
\end{table}
\newpage

\chapter{TEAM MEMBERS AND DIVIDED ROLES}
\begin{table}[h!]
\centering
\caption{Team Member Responsibilities}
\begin{tabularx}{\textwidth}{|p{0.08\textwidth}|p{0.18\textwidth}|X|}
\hline
\textbf{S.N.} & \textbf{Member} & \textbf{Primary Responsibilities} \\ \hline
1 & Member 1 & Project Lead, System Analysis, Backend API (FastAPI), JWT/RBAC, Documentation \\ \hline
2 & Member 2 & Frontend Development (React.js, TypeScript), Dashboard Design, UI/UX \\ \hline
3 & Member 3 & Database Design (PostgreSQL), Assignment Module, Analytics \& GPA Prediction \\ \hline
4 & Member 4 & AI Attendance Module (YOLOv8), WebSocket Integration, Docker Deployment \\ \hline
\end{tabularx}
\end{table}

\section{Shared Responsibilities}
All team members will collaboratively contribute to requirement gathering, literature review, testing and debugging, report writing, presentations, deployment, and user acceptance testing.

\begin{thebibliography}{99}
\bibitem{budi2020} Budi, H., et al. (2020). ``IBAtS: Image Based Attendance System Using HAAR Cascade and LBPH.'' \textit{International Journal of Advanced Computer Science and Applications}, 11(3), pp. 45--52.
\bibitem{pratama2024} Pratama, B., et al. (2024). ``Real-Time Face Recognition Attendance System Using YOLOv8 and FaceNet.'' \textit{Proceedings of ICCVAI 2024}, pp. 112--120.
\bibitem{gomes2019} Gomes, G., et al. (2019). ``Class Attendance Management System Using Facial Recognition.'' \textit{International Journal of Engineering Research and Technology}, 8(6), pp. 234--241.
\bibitem{yu2022} Yu, Y., and Wang, W. (2022). ``College Student Management System Based on K-Means Clustering.'' \textit{Journal of Educational Information Technology}, 15(4), pp. 78--89.
\bibitem{subbiah2021} Subbiah, R., et al. (2021). ``Development of a Web-Based Student Database Management System.'' \textit{International Journal of Scientific Research in CSE}, 9(2), pp. 34--42.
\bibitem{bucko2021} Bucko, A., et al. (2021). ``Enhancing JWT Authentication Based on User Behavior.'' \textit{International Journal of Information Security}, 15(3), pp. 56--71.
\bibitem{ananda2024} Ananda, M., et al. (2024). ``Enhancing Low-Resolution Facial Recognition Using YOLOv8.'' \textit{Proceedings of ICAICV 2024}, pp. 88--97.
\bibitem{kle2024} KLE Technological University. (2024). ``YOLOv8-Enhanced One-Shot Learning Attendance System.'' \textit{International Journal of Computer Vision Applications}, 12(1), pp. 23--35.
\bibitem{roy2012} Roy, S., et al. (2012). ``Automated Attendance System Using Face Recognition.'' NIT Agartala Technical Report, NITA-CS-2012-05.
\bibitem{ahmed2024} Ahmed, A., et al. (2024). ``Fine-Tuning YOLOv8s for Crowded Face Detection.'' \textit{International Journal of Computer Vision}, 7(2), pp. 101--115.
\bibitem{smith2024} Smith, D., and Brown, J. (2024). ``Modular Monoliths in Practice.'' \textit{Software Architecture Journal}, 18(1), pp. 45--59.
\bibitem{kumar2024} Kumar, R., et al. (2024). ``Face Recognition Using FaceNet and YOLOv8.'' \textit{International Journal of AI Research}, 9(1), pp. 67--82.
\bibitem{chen2023} Chen, L., et al. (2023). ``Smart Attendance System Using Improved Facial Recognition.'' \textit{Proceedings of ISA 2023}, pp. 45--53.
\bibitem{hub2024} Smart Student Hub Team. (2024). ``Smart Student Hub: Unified Academic Portal.'' \textit{Journal of Web Engineering}, 6(2), pp. 33--48.
\bibitem{dmc2023} DMC College Research Group. (2023). ``Automated Academic Record Management with BI.'' \textit{International Journal of Educational IS}, 11(4), pp. 56--70.
\bibitem{sharma2021} Sharma, A. (2021). ``Student Information System: Design Framework.'' \textit{International Journal of Information Management}, 14(2), pp. 22--38.
\bibitem{patel2023} Patel, K., et al. (2023). ``Student Management System Using Microservices.'' \textit{International Journal of Cloud Computing}, 8(3), pp. 44--58.
\bibitem{jocher2023} Jocher, G., et al. (2023). ``Ultralytics YOLOv8.'' GitHub Repository, Ultralytics Inc.
\bibitem{schroff2015} Schroff, F., et al. (2015). ``FaceNet: A Unified Embedding for Face Recognition.'' \textit{Proceedings of CVPR 2015}, pp. 815--823.
\bibitem{fastapi2024} FastAPI Documentation. (2024). Available: \url{https://fastapi.tiangolo.com/}
\bibitem{react2024} React Documentation. (2024). Available: \url{https://react.dev/}
\bibitem{postgres2024} PostgreSQL Documentation. (2024). Available: \url{https://www.postgresql.org/docs/}
\bibitem{ultralytics2024} Ultralytics Documentation. (2024). Available: \url{https://docs.ultralytics.com/}
\bibitem{docker2024} Docker Documentation. (2024). Available: \url{https://docs.docker.com/}
\bibitem{nginx2024} Nginx Documentation. (2024). Available: \url{https://nginx.org/en/docs/}
\bibitem{mdn2024} MDN WebSocket API. (2024). Available: \url{https://developer.mozilla.org/en-US/docs/Web/API/WebSockets_API}
\bibitem{typescript2024} TypeScript Documentation. (2024). Available: \url{https://www.typescriptlang.org/docs/}
\end{thebibliography}

\end{document}
